\section{Совместный сертификат по зазору и норме}
\label{joint:section}

\begin{theorem}[Точная совместная огибающая]
\label{joint:main}
При условиях \eqref{boundarygap:assumptions} для $g,P>0$ положим
\begin{equation}\label{joint:definition}
 h=\min\{g,P^2/m\},\qquad
 \mathcal J(g,P)=\max_{h^2/P^2\le\rho\le h/m}V(h,\rho),
\end{equation}
а на координатных осях положим $\mathcal J=0$. Тогда
\begin{equation}\label{joint:bound}
 \delta(x)\le\mathcal J\bigl(G(x),\|\nabla f(x)\|\bigr)
 \le\min\{\mathcal B(G(x)),\mathcal E(\|\nabla f(x)\|)\}.
\end{equation}
При $g\ge P^2/m$ имеем $\mathcal J(g,P)=\mathcal E(P)$.
Функция $\mathcal J$ конечна, непрерывна и не убывает по каждому
аргументу на $[0,\infty)^2$.

Для фиксированных $g,P>0$ значение $\mathcal J(g,P)$ равно
супремуму ошибки при точных наблюдениях $G(x)=g$,
$\|\nabla f(x)\|=P$. Тот же супремум получается при верхних
границах $G(x)\le g$, $\|\nabla f(x)\|\le P$.
Для точной пары достаточно гладких примеров в размерности не выше
двух, для верхних границ достаточно отрезков. Области выпуклы,
компактны и имеют непустую внутренность в своём объемлющем
пространстве, но могут меняться. Точность для одной фиксированной
области не утверждается; формула \eqref{joint:definition}
не использует границу диаметра.
\end{theorem}

\begin{proof}\label{joint:proof}
Сначала пусть $G(x)\le g$, $\|\nabla f(x)\|\le P$.
При $x\ne x^*$ радиус $r=\|x-x^*\|$ и конечный радиальный
наклон $p$ удовлетворяют
\[
 p\le\min\{P,g/r\},\qquad
 r\le\min\{P/m,\sqrt{g/m}\}.
\]
При $g\ge P^2/m$ применима нормовая граница.
При $g<P^2/m$ положим $r_0=g/P$.
Для $r\le r_0$ монотонность радиальной огибающей по длине даёт
\[
 \delta(x)\le A(P,r)\le A(P,r_0)=V(g,g^2/P^2).
\]
Для $r\ge r_0$ имеем $\delta(x)\le V(g,r^2)$ и $r^2\le g/m$.
Это доказывает верхнюю границу в \eqref{joint:definition}.
При нулевом наблюдении или $x=x^*$ ошибка равна нулю.

Докажем точность при точной паре. Если $g<P^2/m$, выберем
максимизирующий $\rho$ и положим
$r=\sqrt\rho$, $p=g/r\le P$, $q=p-mr$.
По радиальной лемме существуют гладкие $u_i$ на $[0,r]$ с
\[
 u_i(0)=0,\quad u_i'(0)=q,\quad u_i'(r)=p,\quad
 u_i'\ge q+ms,\quad |u_i''|\le H_0+H_1u_i,
 \quad u_i(r)\longrightarrow V(g,\rho).
\]
При любом $a>0$ положим
\[
 c=\sqrt{P^2-p^2},\qquad Q=[0,r]\times[0,a],\qquad
 f_i(s,y)=u_i(s)+cy+my^2/2.
\]
Минимум единственен и равен $(0,0)$. Сдвинутое RSI следует из
$(u_i'(s)-q)s+my^2\ge m(s^2+y^2)$, а диагональный гессиан даёт
\[
 \|\nabla^2f_i(s,y)\|
 =\max\{|u_i''(s)|,m\}\le H_0+H_1u_i(s)\le H_0+H_1f_i(s,y).
\]
В точке $(r,0)$ градиент равен $(p,c)$, его норма равна $P$,
а полный зазор равен $rp=g$.

Если $g\ge P^2/m$, возьмём стационарные почти экстремальные
$u_i$ на $[0,r]$, $r=P/m$, и продолжим их влево до
$-a$, где $a=g/P-r$. По конструкции сглаживания около нуля
они совпадают с
\[
 u_+(s)=
 \begin{cases}
 (H_0/H_1)\{\cosh(\sqrt{H_1}s)-1\},&H_1>0,\\
 H_0s^2/2,&H_1=0.
 \end{cases}
\]
Та же чётная функция задаёт гладкое продолжение.
На $[-a,0]$ имеем $u_+'(s)s\ge H_0s^2\ge ms^2$ и
$|u_+''|=H_0+H_1u_+$. Поэтому на $Q=[-a,r]$ все условия
сохраняются, норма в $r$ равна $P$, полный зазор равен
$(r+a)P=g$, а ошибка стремится к $\mathcal E(P)$.

Для верхних границ в первом случае уберём вторую координату,
а во втором оставим стационарный отрезок без продолжения.
Получаются одномерные примеры с теми же предельными ошибками.
Отсюда следует второе неравенство в \eqref{joint:bound} и
монотонность по границам, поскольку больший допустимый класс
имеет не меньший супремум. На положительном квадранте непрерывность
следует из непрерывности $V$ и концов отрезка максимизации.
На осях она следует из \eqref{joint:bound}.
При $g=0$ возможна любая точная норма в минимуме; при $P=0<g$
точная пара невозможна, и $\mathcal J(g,0)=0$ относится к
верхним границам.
\end{proof}

\begin{corollary}[Количественное улучшение нормовой границы]
\label{joint:strict}
При $H_1\ge0$ и $0<g<P^2/m$
\[
 \mathcal J(g,P)\le\mathcal E(P)
   -\frac{(P^2-mg)^2}{2mP^2}<\mathcal E(P).
\]
Если также $H_1>0$ и $\mathcal S(g)\le0$, то
$\mathcal J(g,P)<\min\{\mathcal B(g),\mathcal E(P)\}$.
\end{corollary}
\begin{proof}
Для $\rho\in[g^2/P^2,g/m]$ положим
$p=g/\sqrt\rho$, $q=p-m\sqrt\rho$.
Радиальная поправка и рост $\mathcal E$ дают
\[
 \mathcal E(P)-V(g,\rho)
 \ge\mathcal E(P)-\mathcal E(p)+\mathcal E(q)
 \ge\frac{P^2-p^2+q^2}{2m}
 =\frac{P^2-2mg+m^2\rho}{2m}
 \ge\frac{(P^2-mg)^2}{2mP^2}.
\]
Максимизируем по $\rho$. При $\mathcal S(g)\le0$ непрерывная
функция $V(g,\rho)$ строго меньше $g=\mathcal B(g)$ на компактном
отрезке с положительным левым концом.
\end{proof}
Граница разницы достигается при $H_0=m$, $H_1=0$.

\begin{corollary}[Один корень или алгебраическая граница]
\label{joint:ordinary}\label{joint:regime}
Пусть $H_1>0$, $g,P>0$, $h=\min\{g,P^2/m\}$.
Если $\mathcal S(h)\le0$, то
\begin{equation}\label{joint:ordinary-formula}
 \mathcal J(g,P)=V(h,h^2/P^2)
 \le h+\frac{h^2}{P^2}\mathcal S(h)\le h.
\end{equation}
\end{corollary}
\begin{proof}
По следствию~\ref{gapregime:radius-monotone} максимум достигается
в левом конце допустимого отрезка; оценку даёт
\eqref{gapregime:tangent-upper}.
\end{proof}
Для $\mathcal J$ нужен один корень переключения, без внешней сетки.
Правая часть вычисляется арифметикой и одним квадратным корнем.
Эта алгебраическая граница не обещает заданной относительной
точности приближения $\mathcal J$.
Так, при $m=H_0=H_1=1$, $g=1$, $P=2$
\[
 \mathcal S(1)=4-\frac{7\sqrt{14}}6<0,\qquad
 \mathcal J(1,2)\le2-\frac{7\sqrt{14}}{24}<1.
\]
Упрощение применимо лишь при доказанном знаке $\mathcal S(h)\le0$.

\begin{corollary}[Постоянная граница кривизны]
\label{joint:constant}
При $H_1=0$, $g,P>0$, $h=\min\{g,P^2/m\}$ и
$\beta=(H_0^2-2mH_0-m^2)/(4H_0)$
\begin{equation}\label{joint:constant-formula}
 \mathcal J(g,P)=
 \begin{cases}
 h+\beta h/m,&\beta\ge0,\\
 h+\beta h^2/P^2,&\beta<0.
 \end{cases}
\end{equation}
При $m\le H_0<(1+\sqrt2)m$ и $0<g<P^2/m$ совместная оценка
строго меньше обеих отдельных.
\end{corollary}
\begin{proof}
Максимизируем аффинную функцию $V(h,\rho)=h+\beta\rho$.
При $\beta<0$ положим $b=-\beta/m\in(0,1/2]$,
$t=mg/P^2\in(0,1)$. Тогда $\mathcal B(g)=g$ и
\[
 \mathcal J(g,P)=\frac{P^2}{m}(t-bt^2),\quad
 \mathcal E(P)=\frac{P^2}{m}(1-b),\quad
 1-b-t+bt^2=(1-t)\{1-b(1+t)\}>0.
\]
\end{proof}
При $m=H_0=P=1$, $H_1=0$, $g=1/2$ получаем
$\mathcal J=3/8$, тогда как $\mathcal B=\mathcal E=1/2$.
Функция $f(s)=s/2+s^2/2$ на $[0,1/2]$ в правом конце
достигает равенства при точной паре наблюдений.

\paragraph{Общее скалярное вычисление.}
Вне режима \eqref{joint:ordinary-formula} применим
предложение~\ref{joint:budget} с зазором $h$, отрезком
$[h^2/P^2,h/m]$ и $\ell=h/2$.
Оно даёт $\mathcal J\le\widehat{\mathcal J}\le\vartheta\mathcal J$;
при $g\ge P^2/m$ отрезок состоит из одного узла.
Для $H_1=0$ и на координатных осях используются явные формулы.

\paragraph{Неточные наблюдения.}
Из $G(x)\le B^+$ и $\|\nabla f(x)\|\le P^+$ следует
$\delta(x)\le\mathcal J(B^+,P^+)$.
Гарантия $\|\widehat g-\nabla f(x)\|\le\eta_{\rm norm}$ даёт
$P^+=\|\widehat g\|+\eta_{\rm norm}$.
Одна граница направленных ошибок $\eta_{\rm dir}$ или ошибок
разностей оценок атомов такой гарантии нормы не даёт.
Если оба наблюдения уже известны, новых запросов функции,
градиента или LMO не нужно; вычисление нормы и скалярного
сертификата имеет арифметическую стоимость.
Значение $f(x)-\mathcal J(B^+,P^+)$ есть нижняя граница оптимума.
Неравенство $\mathcal J\ge W(h)$ относится к огибающей и не
является нижней границей действительной ошибки $\delta(x)$.
